% =============================================================================
% Nota de projeto para a primeira reunião de orientação — 18/08/2026
% UFOPA / IEG — Bacharelado em Ciência da Computação
%
% Não é o pré-projeto. É o documento curto que abre a conversa: o problema, a
% proposta, o que já foi verificado e o que precisa ser decidido com a
% orientadora. O pré-projeto nasce depois, com o recorte já acertado.
%
% Formatação vem inteira de ../comum/ufopa-tcc.sty. Não configure nada aqui.
% =============================================================================

\documentclass[12pt,a4paper,oneside]{article}

\usepackage{ufopa-tcc}

% O preâmbulo do TCC abre nova página a cada seção primária, como manda a norma
% para monografia. Aqui são duas páginas corridas: restauro o \section original,
% que o próprio pacote guardou em \ufopaoriginalsection.
\let\section\ufopaoriginalsection

% Com seção em página nova, espaço zero acima do título é o certo. Em texto
% corrido, cola o título no parágrafo anterior — daí o respiro antes e a
% redução depois.
\titlespacing*{\section}{0pt}{0.8\baselineskip}{0.2\baselineskip}

% Sem cabeçalho nem número de página: o documento não é peça normalizada.
\pagestyle{empty}

\ufopatitulo{Sistemas inteligentes aplicados à gestão pública}
\ufopasubtitulo{automação da fase preparatória e pesquisa de preços em
  licitações municipais}
\ufopaautor{Rafael Nobre de Souza}
\ufopaorientador{Prof.\textsuperscript{a} Socorro Vânia Lourenço Alves}
\ufopainstituicao{Universidade Federal do Oeste do Pará}
\ufopalocal{Santarém --- Pará}
\ufopaano{2026}

\begin{document}

\begin{center}
  \begin{singlespace}
    {\bfseries\MakeUppercase{Sistemas inteligentes aplicados à gestão pública:}}\\
    {\bfseries automação da fase preparatória e pesquisa de preços em licitações
     municipais}\\[8pt]
    {\small Rafael Nobre de Souza \quad$\vert$\quad
     Orientadora: \ufopaimprimirorientador\\
     Nota de projeto para a primeira reunião de orientação --- 18 de agosto de 2026}
  \end{singlespace}
\end{center}

\section*{1\quad O problema}

A Lei nº 14.133/2021 caracteriza a fase preparatória da licitação como
planejamento e exige que a necessidade seja fundamentada em estudo técnico
preliminar, que o objeto seja definido em termo de referência e que o orçamento
estimado traga as composições dos preços utilizados (art.~18, incisos I, II e
IV). Em um município pequeno isso recai sobre poucos servidores, sem apoio
técnico dedicado, que produzem cada documento manualmente e refazem o trabalho
quando o processo volta com apontamento do controle interno ou do
assessoramento jurídico.

A etapa mais custosa é a pesquisa de preços. A lei quer o valor estimado
compatível com o mercado a partir de bancos de dados públicos, e admite entre
seus parâmetros a mediana do item no painel de preços do Portal Nacional de
Contratações Públicas (PNCP) e as contratações similares concluídas no ano
anterior (art.~23, §1º, incisos I e II). Na prática o servidor recorre à cotação
direta com fornecedores, o parâmetro mais frágil, porque consultar o histórico
de contratações item a item é hoje trabalho manual inviável.

\section*{2\quad A proposta}

Um sistema de apoio à fase preparatória, com dois módulos sobre a mesma base:
\textbf{pesquisa de preços} --- o núcleo do trabalho ---, que monta a cesta de
preços de cada item a partir de contratações públicas já homologadas,
registrando a procedência de cada valor; e \textbf{apoio à redação}, que gera
minutas de estudo técnico preliminar e termo de referência a partir dos dados
já levantados.

O sistema não se conecta ao sistema de compras do município: produz o documento
e a memória de cálculo, que o servidor revisa, altera e assina. A escolha reduz
o risco institucional de adoção e mantém a responsabilidade legal onde a lei a
coloca.

\section*{3\quad Por que isto é um problema de banco de dados}

Verifiquei a API pública do PNCP em 17 e 18 de agosto, com requisições reais.
O preço unitário homologado por item é público, e a varredura por município e
período funciona --- um mês inteiro do Pará traz 821 contratações em cerca de
8 segundos. Mas \textbf{nenhuma das interfaces pesquisa pela descrição do
item}: a busca do portal casa o termo contra o objeto da contratação. Uma
contratação com 223 itens, cujo item 181 é um microscópio digital, não aparece
na busca por \enquote{microscópio digital} no Pará.

A consequência define o trabalho: é preciso percorrer as contratações, baixar
os itens de cada uma e \emph{construir um índice próprio sobre as descrições}
--- texto livre de até 2048 caracteres, escrito por servidor diferente em cada
município, sem catálogo preenchido. \enquote{Papel A4 75~g/m²} e \enquote{resma
de papel sulfite A4 gramatura 75} são o mesmo produto e não casam por igualdade
de texto. Modelagem, indexação e busca por similaridade não são acessório
tecnológico: são a condição do módulo.

Os riscos também foram medidos, e não supostos: a coleta devolve 10 itens de
uma contratação que tem 223 sem sinalizar erro; filtro escrito errado é
ignorado em silêncio; o serviço ficou fora do ar por cerca de 1h40 e falhou de
novo no dia seguinte; um endereço da API mudou de lugar entre 17 e 18 de
agosto, sem aviso. Caracterizar essas limitações é parte do resultado.

\section*{4\quad Como pretendo avaliar}

Duas medidas, sobre uso real em uma prefeitura de pequeno porte: o
\textbf{tempo por etapa}, contra uma linha de base medida antes da adoção; e a
\textbf{completude formal} do documento gerado, verificada automaticamente
contra os elementos dos arts.~6º e 18. É estudo de caso: um único município.

\section*{5\quad O que preciso decidir com a senhora}

\begin{enumerate}\setlength{\itemsep}{0pt}
  \item Os dois módulos cabem em dois semestres, ou o trabalho deve se limitar
    à pesquisa de preços?
  \item O eixo de banco de dados basta, ou prefere puxar para um lado ---
    similaridade, modelagem temporal de preços, auditabilidade?
  \item A contribuição precisa ser uma pergunta de pesquisa com resposta
    medida, ou um sistema avaliado em campo já atende? O desenho da seção 4
    serve?
  \item A Instrução Normativa SEGES/ME nº 65/2021, que regula a pesquisa de
    preços, é federal --- mas o art.~187 permite que municípios apliquem os
    regulamentos da União. Adoto a IN 65 como referência metodológica?
  \item Titulação e natureza do trabalho para a folha de rosto do TCC~I; datas
    de 2026.2; e se o produto pode ser artigo publicado.
\end{enumerate}

% Corpo menor só aqui: esta nota não é peça normalizada, e as duas referências
% legislativas têm URL longa. No TCC vale o tamanho do pacote, sem este ajuste.
\renewcommand*{\bibfont}{\small\singlespacing}

\nocite{brasil2021lei14133,brasil2021in65}
\printbibliography[title={Referências}]

\end{document}
